\documentclass[11pt]{article}
\usepackage[margin=1in]{geometry}
\usepackage{amsmath,amssymb,amsfonts,mathtools}
\usepackage{array,booktabs,enumitem,graphicx,xcolor,hyperref}
\usepackage[T1]{fontenc}
\usepackage[utf8]{inputenc}
\title{On quantum group symmetries of conformal field theories}
\author{Source-backed arXiv sample}
\date{}
\begin{document}
\maketitle
\begin{abstract}
The appearance of quantum groups in conformal field theories is traced back to the Poisson-Lie symmetries of the classical chiral theory. A geometric quantization of the classical theory deforms the Poisson-Lie symmetries to the quantum group ones. This elucidates the fundamental role of chiral symmetries that quantum groups play in conformal models. As a byproduct, one obtains a more geometric approach to the representation theory of quantum groups.
\end{abstract}
\section{References And Citations}
\label{sec:refs}
This sample cites source keys \cite{10} and \cite{11}. Section~\ref{sec:refs} and Equation~\ref{eq:source-demo} provide cross-reference coverage.
\begin{equation}
\label{eq:source-demo}
a^2 + b^2 = c^2
\end{equation}
\begin{thebibliography}{9}
\bibitem{10} Source bibliography entry 1 extracted from arXiv metadata/source key `10`.
\bibitem{11} Source bibliography entry 2 extracted from arXiv metadata/source key `11`.
\bibitem{12} Source bibliography entry 3 extracted from arXiv metadata/source key `12`.
\bibitem{13} Source bibliography entry 4 extracted from arXiv metadata/source key `13`.
\end{thebibliography}

\end{document}
